\section*{Scope and research status}
\addcontentsline{toc}{section}{Scope and research status}
This paper treats the supplied Alp\"oge and OpenAI theorems as cited
results, as requested. It distinguishes those premises from deductions
proved in this run. ``Finding'' means such a deduction, including a
counterexample or an elementary sharpened form of a known result; it does
not assert historical priority. Exact source locators accompany the
statements used. The accompanying check register describes what was
actually checked and does not promote a text scan or a numerical test to
a formal proof certificate.

A central positive result obtained here is P-F11: a global
quasi-polynomial with nonnegative rational Hilbert numerator and
$r$ denominator factors of degrees at most $A$ has positivity index
less than $rA+128r^3A^3$. P-F12 applies this to show that the revised
PH3 of GCT VI implies its revised PH2. This is an implication from
bounded positive Hilbert data, not a proof that the requested
denominator bounds exist.
P-F14--P-F15 separately refute the rank-only PH3, PH2 and saturation
bounds in the broader subgroup-restriction formulation: faithful
polynomial torus embeddings have fixed group ranks and unbounded
indices as their character weights vary. The narrower plethysm
formulation is not refuted by those examples.
P-F16--P-F18 prove finite uniform bounds for a fixed semisimple
embedding and, for ordinary $GL_n$ plethysm, for fixed inner dimension.
P-F19 then gives an exponential lower bound for that optimal canonical
modular-index envelope, using the single inner shape $(3)$.
It refutes the stronger canonical clause of ordinary plethysm PH3.
The weaker assertion about some positive rational Hilbert expression,
and the narrower PH2 and saturation bounds, remain separate questions.
P-F20 distinguishes a large universal denominator from an ordinary
Hilbert-series pole. P-F21--P-F22 refute general transfers from large
parameter degrees even for rational Gorenstein domains with one fixed
Hilbert series. P-F23 supplies a precise noncancellation criterion;
P-F24 proves that the chain gives a smooth one-dimensional fixed
component in the actual cubic quotient. Controlling the other fixed
components is the remaining geometric gate on this route to the
weaker positive-form clause. P-F26 confines every competing component
to the discriminant boundary. P-F27 proves that any such component
requires three-distinct-variable monomials in every finite-lift
eigenbasis; even a connected graph part forces the smooth chain.
P-F29 excludes the complementary cube-and-squarefree support class,
so the remaining supports mix the two types with disconnected
graph part. P-F28 gives the exact local chain coefficient and its
cyclotomic valuation, without a global noncancellation assertion.
Positive Hilbert forms also yield
shifted polytope models (P-F2), and the main counterexample ray admits
an optimal two-dimensional model (P-F1).

The converse direction from qualitative purity to the required counting
and complexity properties fails in concrete examples. PLE-F7 supplies a
normal surface with pure even Tate ordinary and intersection cohomology
whose Hilbert function violates homogeneous PH1. PLE-F8 identifies it
with the actual canonical plethysm multiplicity algebra on that ray,
including multiplication, and supplies a basis with nonnegative
multiplication constants. P-F10 keeps the smooth
underlying variety fixed while its grading indices grow arbitrarily.
P-F7 shows that finite smooth proper pure Tate schemes can encode all
\#3SAT counts when specified by equations. These examples identify
missing hypotheses; they do not prove a logical independence theorem
between RH and any GCT assertion.

No separation or equality of $\mathsf P$ and $\mathsf{NP}$, or of the
usual complex algebraic classes $\mathsf{VP}$ and $\mathsf{VNP}$, is
proved. The remaining research targets are displayed separately in
Part~II and in \texttt{QUESTIONS.tsv}. A statement is not marked proved
merely because its missing input has been named.

References to numbered GCT VI hypotheses concern the supplied
arXiv:0704.0229v4 text. The broad subgroup-restriction formulation
of Problem 1.1.3 and Sections 3.3--3.4 is distinguished throughout
from irreducible-inner Problem 1.1.2 and its ordinary $GL_n$
plethysm subfamily. Introductory notational inconsistencies about
the canonical invariant algebra are resolved using the explicit
construction and modular-index definition in Sections 4.1.1--4.1.2.

\subsection*{The supplied analytic statements}
The five theorem declarations in \texttt{ALPOGE\_Main.lean} are the
precise analytic premises. Their types are imported from the supplied
\texttt{ALPOGE\_Statements.lean}, except for Mathlib's
\texttt{RiemannHypothesis} itself.
\begin{longtable}{@{}p{.48\linewidth}p{.46\linewidth}@{}}
\toprule
Theorem name (statement identifier) & Scope used here\\
\midrule
\endhead
\path{riemannHypothesis} (A-RH) & Mathlib's named Riemann-hypothesis proposition.\\[3pt]
\path{riemannZeta_eq_zero_re} (A-ZETA) & A zero of $\zeta(s)$, other than a negative even integer and with $s\ne1$, has $\Re s=1/2$.\\[3pt]
\path{LFunction_ne_zero_of_half_lt_re} (A-DIR-NV) & For every positive modulus and Dirichlet character, $L(\chi,s)\ne0$ if $\Re s>1/2$ and $(\chi\ne1\text{ or }s\ne1)$.\\[3pt]
\path{LFunction_eq_zero_re} (A-DIR-Z) & All Dirichlet $L$-zeros in $0<\Re s<1$ have real part $1/2$.\\[3pt]
\path{LFunction_eq_zero_re_of_isPrimitive} (A-DIR-PRIM) & The corresponding statement restricted to primitive Dirichlet characters.\\
\bottomrule
\end{longtable}
These declarations quantify over Dirichlet characters, not arbitrary
number fields or their Hecke characters. C-GRH-01 proves the exact
cyclotomic-field extension needed below. The present working set contains
the theorem wrappers, their statement file and a supplied scan report;
it does not contain all imported proof modules. No Lean kernel rebuild
was performed. This records the scope of this run's verification, without
reopening the requested treatment of the theorems as cited results.

\subsection*{The positivity method in the supplied note}
The determinant sign $W(1_S)\ge0$ is in \emph{Remark}, p.~3, numbered
item~4. The sign-only discussion of finite-field purity is in footnote~3
on that page; numbered item~5 concerns switching to the von Mangoldt
sum. GGP periods and interlacing appear in item~7, iterated theta lifts in
item~8, and the Bessel-period seesaw in item~12 on p.~4. We use these
locators rather than conflating the footnote with numbered item~5.

The transfer question is mathematical: identify the representation
multiplicity or graded coordinate-ring quotient with the object to which
a sign or purity theorem applies, then supply quantitative control.
P-F3 constructs a positive determinant for each coefficient
unconditionally, exposing the dimension cost. P-F4 separates pointwise
sign from character positivity. P-F5 distinguishes standard interlacing
from restriction along a Schur-functor homomorphism. P-F6 states the
precise categorical data sufficient for positivity of a based product.
Neither the four-page note nor the five theorem types supplies those GCT
identifications automatically.

\subsection*{Reading the package}
Part~I contains the checked findings and the exact cited background on
which they rely. Part~II contains unresolved targets and unsuccessful
transfer attempts; it asserts no unchecked theorem. The source archive
is identified by a SHA-256 manifest. The audit register records model,
effort disclosure, agent/run, date, method and outcome for each check.
The finite computations are sanity checks, separately labelled from
whole-proof readings and independent re-proofs.
